\section{Contexto macro: por que condicionar, e não empilhar}
\label{sec:debate-macro}

\begin{confronto}
``Adicionar Selic, IPCA, dólar e VIX ao tensor de entrada é o caminho óbvio para dar contexto
sistêmico. Por que inventar um \emph{Conditional} Autoencoder em vez de simplesmente passar de
$(30,2)$ para $(30,2{+}N)$ e reconstruir tudo?''
\end{confronto}

Porque ``empilhar e reconstruir tudo'' \textbf{não atinge o objetivo} --- e isso é demonstrável,
não opinião. Três falhas:

\begin{enumerate}[noitemsep]
  \item \textbf{O MSE ignora a macro.} Indicadores de baixa frequência (IPCA mensal, Selic em
        escada), após o preenchimento causal, são quase constantes na janela. Sua variância escalada
        é ínfima; reconstruí-los é trivial; o gradiente vem só de preço/volume. A macro entraria como
        \emph{decoração inerte}.
  \item \textbf{Eventos agendados viram falsos positivos.} Uma alta de Selic ou um \emph{print} de
        IPCA são \emph{programados}. Se a macro está no alvo e o erro é agregado por \emph{max}
        (Seção~\ref{sec:debate-fase2-pooling}), um evento de calendário dispararia uma ``anomalia''.
  \item \textbf{Um escalar único apaga a distinção.} Colapsar preço/volume e macro num só erro
        impede dizer \emph{qual} bloco se mexeu --- justamente a informação que separa idiossincrático
        de sistêmico.
\end{enumerate}

A solução é o \textbf{Conditional Autoencoder}: o encoder \emph{vê} a macro (ela condiciona a
codificação), mas o decoder reconstrói \textbf{só} preço/volume e a perda recai só sobre esse bloco.
O estresse macro é medido \emph{direto da entrada} (amplitude na janela), em um eixo separado do
erro do ativo.

\begin{tradeoff}
Ganha-se a distinção idiossincrático \emph{vs.} sistêmico ao custo de uma arquitetura assimétrica
(entrada $\ne$ alvo) e de um conector de dados macro com tratamento de frequências mistas. A prova
de conceito (notebook \texttt{12}) pagou: COVID/2020 saiu \textbf{sistêmico}; Americanas/2023,
\textbf{idiossincrático}.
\end{tradeoff}

\begin{honesto}
A defesa anti-vazamento depende de um \textbf{contrato}: a macro é indexada pela \emph{data de
publicação}, não de referência (o IPCA de março só existe para a rede em meados de abril), com
\emph{ffill} causal e sem \emph{backfill}. Assumimos um \emph{lag} de publicação \textbf{aproximado}
para o IPCA --- não a data exata de divulgação ---, o que é uma simplificação a refinar. E as séries
mensais (Selic, IPCA) acabam quase inertes na janela: o contexto sistêmico vem, na prática, dos
indicadores diários (USD/BRL, VIX). O método funciona, mas o crédito honesto vai para os dois canais
diários, não para os quatro.
\end{honesto}

\begin{honesto}
Esta extensão (ADR-0012) é uma \textbf{prova de conceito} --- ainda não incorporada como entrega
oficial. O \texttt{latent\_dim} do espaço ampliado carece de revalidação por \emph{walk-forward}
antes de qualquer número final.
\end{honesto}
